\chapter{Asignación de pines}
\label{anx:pines}

\section{Raspberry Pi Pico 2 (revisión B de la placa)}

\begin{longtable}{@{}lllL{0.44\textwidth}@{}}
\caption{Asignación de GPIO según el esquema de la revisión B y el firmware.}\label{tab:pines}\\
\toprule
GPIO & Pin & Señal & Función \\
\midrule
\endfirsthead
\toprule
GPIO & Pin & Señal & Función \\
\midrule
\endhead
GP0  & 1  & INT1    & Interrupción de la célula 1 (MD1) \\
GP1  & 2  & EN1     & Habilitación / reinicio de las placas Click (\codigo{LC1\_EN\_GPIO}) \\
GP2  & 4  & SDA1    & I\textsuperscript{2}C1 SDA (\codigo{LC\_SDA\_PIN}) \\
GP3  & 5  & SCL1    & I\textsuperscript{2}C1 SCL (\codigo{LC\_SCL\_PIN}) \\
GP4  & 6  & INT2    & Interrupción de la célula 2 (MD2) \\
GP5, GP6 & 7, 9 & -- & Libres \\
GP7  & 10 & SIN4    & Entrada del canal 4 del TPS4H000 \\
GP8  & 11 & SIN3    & Entrada del canal 3 \\
GP9  & 12 & SIN2    & Entrada del canal 2 \\
GP10 & 14 & SIN1    & Entrada del canal 1 \\
GP11 & 15 & DST     & Fallo global del TPS4H000 (drenador abierto, \emph{pull-up}) \\
GP12 & 16 & DE      & Dirección del transceptor RS-485 (\codigo{MB\_EN\_GPIO}) \\
GP13 & 17 & RO      & Recepción RS-485, UART PIO (\codigo{MB\_RX\_GPIO}) \\
GP14 & 19 & DI      & Transmisión RS-485, UART PIO (\codigo{MB\_TX\_GPIO}) \\
GP15 & 20 & DIAG\_EN & Habilitación del diagnóstico del TPS4H000 \\
GP16 & 21 & DIN1    & Final de carrera A (\codigo{PIN\_LIMIT\_A}) \\
GP17 & 22 & DIN2    & Final de carrera B (\codigo{PIN\_LIMIT\_B}) \\
GP18 & 24 & DIN3    & Entrada digital libre \\
GP19 & 25 & DIN4    & Entrada digital libre \\
GP20 & 26 & COMP    & Salida del comparador (umbral \SI{2.5}{\volt}) \\
GP21, GP22 & 27, 29 & -- & Libres \\
GP26 & 31 & AOUT3   & ADC0: entrada analógica AIN3 \\
GP27 & 32 & AOUT1   & ADC1: entrada analógica AIN1 \\
GP28 & 34 & AOUT2   & ADC2: entrada analógica AIN2 \\
VBUS & 40 & +5V     & Alimentación de los amplificadores \\
3V3  & 36 & +3.3V   & Transceptor RS-485 y placas Click \\
\bottomrule
\end{longtable}

\noindent\textbf{Nota.} El firmware define \codigo{LC1\_INT\_GPIO = 4}, mientras que en el esquema de
la revisión B la interrupción de la célula 1 llega a GP0 y GP4 corresponde a la célula 2. El firmware
sólo usa ese pin para llevarlo a nivel bajo durante el apagado de la célula
(sección~\ref{sec:prob-zsc}) \verificar{comprobar a qué GPIO está conectado INT en la placa montada}.

\section{Bornas de campo}

\begin{table}[H]
\centering
\caption{Bornas de la placa (revisión B).}
\label{tab:bornas}
\begin{tabular}{@{}lll@{}}
\toprule
Borna & Pin & Señal \\
\midrule
J5 (RS-485) & 1 / 2 / 3 & A / B / GND \\
Terminal de salidas & 1--8 & SOUT4, SOUT3, --, VS\_IN, SOUT2, SOUT1, AIN2, AIN1 \\
Terminal de entradas & 1--8 & AIN3, AIN4, VS\_IN, --, DI1, DI2, DI3, DI4 \\
Filas de masa & 1--8 & GND \\
\bottomrule
\end{tabular}
\end{table}

\verificar{comprobar el orden de las bornas sobre la placa fabricada}
